% CODEX: wstaw w 30b-geometrizacja-thurstona.tex bezpośrednio przed
% \textbf{Geometria $\mathrm{Nil}$.}, po istniejącym opisie
% \widetilde{\mathrm{SL}}(2,\R). Nie powtarzaj istniejącego
% stwierdzenia o U\Sigma_g jako zwartym ilorazie.

\begin{przyklad}[Jawna metryka geometrii $\widetilde{\mathrm{SL}}(2,\R)$]
\label{prz:thur-dop-tilde-sl-metryka}
Weźmy model półpłaszczyzny hiperbolicznej
$\mathbb H^2=\{(x,y):y>0\}$ z metryką
$h=y^{-2}(\dd x^2+\dd y^2)$ oraz zorientowaną
ramą ortonormalną $e_1=y\partial_x$, $e_2=y\partial_y$.
Wzory na symbole Christoffela z
Przykładu~\ref{prz:christoffel-hiperboliczna} dają
$\nabla e_1=A\,e_2$ i $\nabla e_2=-A\,e_1$,
gdzie $A=\dd x/y$ jest formą koneksji tej ramy.
Niech $\varphi$ oznacza kąt jednostkowego wektora
$v=\cos\varphi\,e_1+\sin\varphi\,e_2$.
W tych współrzędnych metryka Sasakiego na $U\mathbb H^2$
przyjmuje postać
\begin{equation}\label{eq:thur-dop-tilde-sl-metryka}
 g=\frac{\dd x^2+\dd y^2}{y^2}
    +\left(\dd\varphi+\frac{\dd x}{y}\right)^2.
\end{equation}
Na wiązce jednostkowych wektorów kąt jest okresowy
($\varphi\in\R/2\pi\Z$); po przejściu do nakrycia
uniwersalnego rozwijamy go do $\varphi\in\R$,
uzyskując jedną z metryk modelu
$\widetilde{\mathrm{SL}}(2,\R)$.

Forma $\alpha=\dd\varphi+\dd x/y$ znika na
wektorach poziomych, a jej krzywizna wynosi
\[
 \dd\alpha=\frac{\dd x\wedge\dd y}{y^2}
     =\mathrm{vol}_{\mathbb H^2}.
\]
Jeżeli krzywa w bazie jest brzegiem obszaru $D$,
to poziome podniesienie zmienia kąt o
$\Delta\varphi=-\int_D\mathrm{vol}_{\mathbb H^2}$
(modulo $2\pi$ w wiązce $U\mathbb H^2$).
W iloczynie $\mathbb H^2\times\R$ mamy natomiast
płaską koneksję produktową $\dd\varphi$.
Właśnie niezerowa krzywizna formy $\alpha$ odróżnia
te dwie geometrie jednorodne. Dalszą konsekwencję
topologiczną obliczymy w
Przykładzie~\ref{prz:thur-dop-seifert-euler}.
\end{przyklad}
